\section{Containerisering}
\label{sec:docker}

VideoAPI distribueres som en Docker-container. Hele driftsmiljøet, inkludert API, metrikkinnsamling og dashboard, orkestreres av en enkelt \texttt{compose.yaml} i prosjektets rot. Målet med oppsettet er at en hvilken som helst Linux-VM med Docker installert skal kunne starte hele stacken med én kommando, samtidig som tjenestene kjører isolert fra hverandre og uten unødvendige privilegier. Dette oppfyller krav I1.1 \cite{hdo_kravspec} om containerisert leveranse.

\subsection{Image-bygging}
\label{subsec:docker-image}

Dockerfilen til VideoAPI er bygd som et flertrinns-image (multi-stage build) for å holde det endelige imaget lite og fritt for byggeverktøy. Første trinn, \texttt{base}, baseres på \texttt{mcr.microsoft.com/dotnet/aspnet:10.0} og legger til \texttt{curl} slik at containerens egen helsesjekk kan polle API-et. Andre trinn, \texttt{publish}, bruker det fullstendige .NET 10 SDK-imaget til å gjenopprette pakker (\texttt{dotnet restore}) og publisere applikasjonen i Release-konfigurasjon. Det ferdige \texttt{final}-trinnet kopierer kun den publiserte utdataen inn i runtime-imaget og setter \texttt{USER \$APP\_UID}, slik at prosessen kjører som en ikke-privilegert bruker. Dette er en anbefalt praksis i offisielle .NET-images for å redusere angrepsflaten \cite{aspnet_rootless}.

Filtreringen skjer via \texttt{.dockerignore}, som ekskluderer \texttt{.git/}, \texttt{.env}, \texttt{bin/}, \texttt{obj/}, lokale IDE-filer og hemmeligheter. Det forhindrer at sensitive verdier eller utvikleroppsett havner i imaget.

\subsection{Compose-stack}
\label{subsec:docker-compose}

Stacken som defineres i \texttt{compose.yaml} består av tre tjenester som alle deler et internt Docker-nettverk og kan referere til hverandre på tjenestenavn. Tabell~\ref{tab:compose-services} viser tjenestene og hvordan de eksponeres mot host.

\begin{table}[h]
\centering
\caption{Tjenester i Compose-stacken for HDO-VideoAPI.}
\label{tab:compose-services}
\begin{tabular}{l l l l}
\toprule
\textbf{Tjeneste} & \textbf{Image} & \textbf{Portmapping} & \textbf{Formål} \\
\midrule
\texttt{hdo-videoapi} & bygd lokalt fra \texttt{Dockerfile} & 80:8080, 443:8443 & API mot omverdenen \\
\texttt{prometheus}   & \texttt{prom/prometheus:latest}     & 9090:9090         & Skraping av metrikker \\
\texttt{grafana}      & \texttt{grafana/grafana:latest}     & 3000:3000         & Visualisering \\
\bottomrule
\end{tabular}
\end{table}

\newpage
API-tjenesten lytter internt på 8080 (HTTP) og 8443 (HTTPS), og videresendes på host til standardportene 80 og 443. TLS-sertifikater monteres skrivebeskyttet fra hostens \texttt{./certs/}-katalog inn til \texttt{/app/certs} i containeren, slik at sertifikater kan rotere uten at bildet må bygges på nytt. Hemmeligheter (API-nøkler, Keycloak-secrets, Grafana-passord) injiseres via \texttt{env\_file: .env} og forblir utenfor bildet og versjonskontrollen.

Prometheus konfigureres av \texttt{prometheus.yml}, som definerer ett skrape-job mot \texttt{hdo-videoapi:8080/metrics} med 15 sekunders intervall. Siden tjenestene deler nettverk, kan Prometheus nå API-et på tjenestenavnet uten DNS-oppsett. Grafana provisjoneres automatisk fra \texttt{./grafana-provisioning/} ved oppstart, slik at dashboards og datakilder gjenopprettes også når volumet tilbakestilles.

\subsection{Tilstandshåndtering og oppstartsrekkefølge}
\label{subsec:docker-state}

Selv om VideoAPI selv er tilstandsløs (jf.\ krav I1.2 \cite{HDOKravspec}), skal
Prometheus' egne tidsserier og Grafanas konfigurasjon overleve omstart av
containerne. Begge tjenestene bruker derfor navngitte Docker-volumer
(\texttt{prometheus\_data} og \texttt{grafana\_data}) som ligger utenfor
containerens skrivelag.

Tjenestene starter i kontrollert rekkefølge slik at metrikkstacken ikke logger
feil før API-et er klart. \texttt{prometheus} har
\texttt{depends\_on: hdo-videoapi} med betingelsen \texttt{service\_healthy}, og
\texttt{grafana} har \texttt{depends\_on: prometheus}. Det betyr at Prometheus
først starter når API-ets \texttt{/health/live}-endepunkt returnerer 200, og at
Grafana først starter etter Prometheus. Selve helsesjekken er definert i
Compose-filen og bruker \texttt{curl} til å treffe \texttt{/health/live} hvert
tiende sekund med fem sekunders timeout og opp til fem nye forsøk før
containeren markeres som \texttt{unhealthy}.

Alle tre tjenestene bruker \texttt{restart: unless-stopped}, slik at de kommer
opp igjen automatisk etter krasj eller VM-omstart, men forblir nede dersom de
er manuelt stoppet av driftspersonell.

\subsection{Ressurskontroll og sikkerhet}
\label{subsec:docker-security}

For å hindre at en enkelt container uttømmer hostens ressurser, settes
\texttt{deploy.resources.limits} per tjeneste: VideoAPI begrenses til 512~MB
minne og én CPU-kjerne, Prometheus til 256~MB, og Grafana til 1~GB. Verdiene er
forankret i den observerte ressursbruken under last- og integrasjonstesting
(kapittel~\ref{ch:testing}), og gir tydelige grenser for systemovervåkningen.

Sikkerhetsmessig hviler stacken på flere lag. API-et selv kjører som
ikke-priviligert bruker (UID 1654 i .NET-imaget), TLS-sertifikatene monteres
skrivebeskyttet, hemmeligheter ligger kun i \texttt{.env}, og
\texttt{appsettings.Docker.json} inneholder kun ikke-hemmelige
produksjonsoverstyringer. Prometheus eksponerer ingen autentisering ut av
boksen og må derfor regnes som halvoffentlig; Grafana beskyttes av eget
brukernavn og passord satt via miljøvariabler.

\subsection{Driftsflyt}
\label{subsec:docker-ops}

Med stacken på plass blir hele drifts-livssyklusen redusert til et fåtall
kommandoer. Hele stacken bygges og startes i bakgrunnen med
\texttt{docker compose up -d --build}, og en enkelttjeneste kan rebygges med
samme kommando etterfulgt av tjenestenavnet. \texttt{docker compose down}
stopper alle containerne, mens \texttt{docker compose down -v} i tillegg
sletter de navngitte volumene. Logger fra en spesifikk tjeneste følges med
\texttt{docker compose logs -f <tjeneste>}. Samspillet mellom helsesjekk,
oppstartsrekkefølge og volumer gjør at en VM kan rebootes uten manuell
inngripen --- alle tjenestene kommer opp i riktig rekkefølge og gjenoppretter
tidligere tilstand.